% !TEX program = pdflatex
% Multi-chapter report template -- BibLaTeX + Biber.
% Build:  latexmk main.tex
% Clean:  latexmk -c
\documentclass[11pt,a4paper]{report}

% --- Encoding & fonts (pdfLaTeX) --------------------------------------------
\usepackage[T1]{fontenc}
\usepackage[utf8]{inputenc}

% --- Core packages ----------------------------------------------------------
\usepackage[margin=1in]{geometry}
\usepackage{microtype}
\usepackage{amsmath}
\usepackage{amssymb}
\usepackage{graphicx}
\usepackage{booktabs}

% --- Bibliography: BibLaTeX with the Biber backend --------------------------
\usepackage[backend=biber,style=numeric,sorting=nyt]{biblatex}
\addbibresource{refs.bib}

% --- Hyperlinks & smart cross-references ------------------------------------
\usepackage{hyperref}
\usepackage{cleveref}

\graphicspath{{figures/}}

\title{A Multi-Chapter Report}
\author{Author One\\ \texttt{author@example.org}}
\date{\today}

\begin{document}
\maketitle

\begin{abstract}
Placeholder abstract for the report. Summarise the scope, approach, and
principal findings in a few sentences.
\end{abstract}

\tableofcontents

% Each chapter lives in its own file under sections/.
% Chapter file -- included from main.tex via % Chapter file -- included from main.tex via % Chapter file -- included from main.tex via \input{sections/intro}.
% Do NOT add \documentclass or \begin{document} here.
\chapter{Introduction}
\label{ch:intro}

Placeholder introduction chapter. This file is pulled into the report with
\verb|\input{sections/intro}|, which keeps the master file short and lets you
work on one chapter at a time.

\section{Background}
\label{sec:background}
Placeholder background section. Cross-reference other parts of the document
with \verb|\Cref{...}| and cite with \verb|\cite{...}|~\cite{placeholder-report}.

\section{Objectives}
\label{sec:objectives}
Placeholder objectives section.
.
% Do NOT add \documentclass or \begin{document} here.
\chapter{Introduction}
\label{ch:intro}

Placeholder introduction chapter. This file is pulled into the report with
\verb|% Chapter file -- included from main.tex via \input{sections/intro}.
% Do NOT add \documentclass or \begin{document} here.
\chapter{Introduction}
\label{ch:intro}

Placeholder introduction chapter. This file is pulled into the report with
\verb|\input{sections/intro}|, which keeps the master file short and lets you
work on one chapter at a time.

\section{Background}
\label{sec:background}
Placeholder background section. Cross-reference other parts of the document
with \verb|\Cref{...}| and cite with \verb|\cite{...}|~\cite{placeholder-report}.

\section{Objectives}
\label{sec:objectives}
Placeholder objectives section.
|, which keeps the master file short and lets you
work on one chapter at a time.

\section{Background}
\label{sec:background}
Placeholder background section. Cross-reference other parts of the document
with \verb|\Cref{...}| and cite with \verb|\cite{...}|~\cite{placeholder-report}.

\section{Objectives}
\label{sec:objectives}
Placeholder objectives section.
.
% Do NOT add \documentclass or \begin{document} here.
\chapter{Introduction}
\label{ch:intro}

Placeholder introduction chapter. This file is pulled into the report with
\verb|% Chapter file -- included from main.tex via % Chapter file -- included from main.tex via \input{sections/intro}.
% Do NOT add \documentclass or \begin{document} here.
\chapter{Introduction}
\label{ch:intro}

Placeholder introduction chapter. This file is pulled into the report with
\verb|\input{sections/intro}|, which keeps the master file short and lets you
work on one chapter at a time.

\section{Background}
\label{sec:background}
Placeholder background section. Cross-reference other parts of the document
with \verb|\Cref{...}| and cite with \verb|\cite{...}|~\cite{placeholder-report}.

\section{Objectives}
\label{sec:objectives}
Placeholder objectives section.
.
% Do NOT add \documentclass or \begin{document} here.
\chapter{Introduction}
\label{ch:intro}

Placeholder introduction chapter. This file is pulled into the report with
\verb|% Chapter file -- included from main.tex via \input{sections/intro}.
% Do NOT add \documentclass or \begin{document} here.
\chapter{Introduction}
\label{ch:intro}

Placeholder introduction chapter. This file is pulled into the report with
\verb|\input{sections/intro}|, which keeps the master file short and lets you
work on one chapter at a time.

\section{Background}
\label{sec:background}
Placeholder background section. Cross-reference other parts of the document
with \verb|\Cref{...}| and cite with \verb|\cite{...}|~\cite{placeholder-report}.

\section{Objectives}
\label{sec:objectives}
Placeholder objectives section.
|, which keeps the master file short and lets you
work on one chapter at a time.

\section{Background}
\label{sec:background}
Placeholder background section. Cross-reference other parts of the document
with \verb|\Cref{...}| and cite with \verb|\cite{...}|~\cite{placeholder-report}.

\section{Objectives}
\label{sec:objectives}
Placeholder objectives section.
|, which keeps the master file short and lets you
work on one chapter at a time.

\section{Background}
\label{sec:background}
Placeholder background section. Cross-reference other parts of the document
with \verb|\Cref{...}| and cite with \verb|\cite{...}|~\cite{placeholder-report}.

\section{Objectives}
\label{sec:objectives}
Placeholder objectives section.


\chapter{Analysis}
\label{ch:analysis}
Placeholder analysis chapter. Refer back to \Cref{ch:intro} as needed and
cite sources such as~\cite{placeholder-report}.

% ----------------------------------------------------------------------------
% Appendices
% ----------------------------------------------------------------------------
\appendix

\chapter{Supplementary Material}
\label{ch:appendix}
Placeholder appendix content.

\printbibliography[heading=bibintoc]

\end{document}
